\chapter{Aljabar: Persamaan dan Pertidaksamaan}\label{ch:g10:algebra}

Aljabar adalah seni berhitung dengan huruf. Bab ini mengulang dan
mempertajam dua langkah dasarnya --- \omterm{def:g10:algebra:expand}{menjabarkan} dan \omterm{def:g10:algebra:expand}{memfaktorkan} --- lalu
memakai keduanya untuk menyelesaikan \omterm{def:g10:algebra:equation}{persamaan} dan pertidaksamaan secara
sistematis, dengan \omterm{def:g10:algebra:signtable}{tabel tanda} sebagai alat utamanya.

\section{Menjabarkan dan memfaktorkan}

\begin{definition}[Menjabarkan, memfaktorkan]\label{def:g10:algebra:expand}
\emph{Menjabarkan}\index{menjabarkan} sebuah bentuk aljabar berarti mengubah
hasil kali menjadi jumlah; \emph{memfaktorkan}\index{memfaktorkan} berarti
sebaliknya, mengubah jumlah menjadi hasil kali. Aturan dasarnya adalah sifat
distributif:
\[
k(a + b) = ka + kb,
\qquad
(a + b)(c + d) = ac + ad + bc + bd .
\]
\end{definition}

\begin{example}\label{ex:g10:algebra:expand}
Jabarkan $(2x + 3)(x - 5)$ selangkah demi selangkah:
\begin{align*}
(2x + 3)(x - 5)
&= 2x \times x + 2x \times (-5) + 3 \times x + 3 \times (-5) \\
&= 2x^2 - 10x + 3x - 15 \\
&= 2x^2 - 7x - 15 .
\end{align*}
\end{example}

\begin{theorem}[Identitas istimewa]\label{thm:g10:algebra:identities}
Untuk semua \omterm{def:g10:numbers:sets}{bilangan real} $a$ dan $b$:
\[
(a+b)^2 = a^2 + 2ab + b^2,
\qquad
(a-b)^2 = a^2 - 2ab + b^2,
\qquad
(a+b)(a-b) = a^2 - b^2 .
\]
\end{theorem}

\begin{proof}
Masing-masing merupakan penjabaran langsung. Untuk yang pertama:
\[
(a+b)^2 = (a+b)(a+b) = a^2 + ab + ba + b^2 = a^2 + 2ab + b^2 .
\]
Yang kedua diperoleh dengan mengganti $b$ oleh $-b$: $(a-b)^2 = a^2 + 2a(-b) +
(-b)^2 = a^2 - 2ab + b^2$. Untuk yang ketiga:
$(a+b)(a-b) = a^2 - ab + ba - b^2 = a^2 - b^2$.
\end{proof}

\begin{omfigure}
\begin{tikzpicture}[scale=0.62]
  \draw[thick] (0,0) rectangle (5,5);
  \draw[thick] (3.4,0) -- (3.4,5);
  \draw[thick] (0,3.4) -- (5,3.4);
  \fill[omDef!12] (0,0) rectangle (3.4,3.4);
  \fill[omThm!12] (3.4,3.4) rectangle (5,5);
  \fill[omProp!12] (3.4,0) rectangle (5,3.4);
  \fill[omProp!12] (0,3.4) rectangle (3.4,5);
  \node[font=\small] at (1.7,1.7) {$a^2$};
  \node[font=\small] at (4.2,4.2) {$b^2$};
  \node[font=\small] at (4.2,1.7) {$ab$};
  \node[font=\small] at (1.7,4.2) {$ab$};
  \node[below, font=\small] at (1.7,0) {$a$};
  \node[below, font=\small] at (4.2,0) {$b$};
  \node[left, font=\small] at (0,1.7) {$a$};
  \node[left, font=\small] at (0,4.2) {$b$};
\end{tikzpicture}

{\small Mengapa $(a+b)^2 = a^2 + 2ab + b^2$: persegi bersisi $a + b$ terbagi
menjadi persegi bersisi $a$, persegi bersisi $b$, dan dua persegi panjang
$a \times b$.}
\end{omfigure}

\begin{example}\label{ex:g10:algebra:identities}
Dipakai maju (\omterm{def:g10:algebra:expand}{menjabarkan}) dan mundur (\omterm{def:g10:algebra:expand}{memfaktorkan}):
\begin{align*}
(3x + 2)^2 &= 9x^2 + 12x + 4,
& x^2 - 25 &= (x+5)(x-5), \\
(x - 4)^2 &= x^2 - 8x + 16,
& 4x^2 + 4x + 1 &= (2x + 1)^2 .
\end{align*}
Berhitung di kepala pun terbantu: $101 \times 99 = (100+1)(100-1) =
10000 - 1 = 9999$.
\end{example}

\begin{method}[Memfaktorkan sebuah bentuk aljabar]\label{met:g10:algebra:factor}
Cobalah, menurut urutan ini:
\begin{enumerate}
  \item \emph{Faktor persekutuan:} temukan faktor yang hadir pada setiap suku
        lalu keluarkan, misalnya $6x^2 + 4x = 2x(3x + 2)$. Faktor persekutuan
        itu boleh berupa satu kurung utuh:
        $(x+1)(2x-3) + (x+1)(x+7) = (x+1)(3x+4)$.
  \item \emph{Identitas istimewa:} kenali $a^2 - b^2$, atau
        $a^2 \pm 2ab + b^2$, misalnya $9x^2 - 16 = (3x)^2 - 4^2 =
        (3x+4)(3x-4)$.
  \item Gabungkan keduanya, lalu periksa hasilnya dengan \omterm{def:g10:algebra:expand}{menjabarkannya} kembali.
\end{enumerate}
\end{method}

\begin{example}\label{ex:g10:algebra:factorsteps}
Faktorkan $E = (2x + 1)^2 - (x - 3)^2$. Ini selisih dua kuadrat
$a^2 - b^2$ dengan $a = 2x+1$ dan $b = x-3$:
\begin{align*}
E &= \bigl[(2x+1) + (x-3)\bigr] \times \bigl[(2x+1) - (x-3)\bigr] \\
  &= (3x - 2)(2x + 1 - x + 3) \\
  &= (3x - 2)(x + 4).
\end{align*}
Perhatikan kurung kedua: mengurangkan $x - 3$ berarti mengurangkan $x$
\emph{dan} menambahkan $3$.
\end{example}

\section{Persamaan}

\begin{definition}[Persamaan, penyelesaian]\label{def:g10:algebra:equation}
Sebuah \emph{persamaan}\index{persamaan} adalah kesamaan yang memuat sebuah
bilangan yang belum diketahui $x$; sebuah \emph{penyelesaian} adalah nilai $x$
yang membuat kesamaan itu benar. Menyelesaikan persamaan berarti mencari
\emph{semua} penyelesaiannya.
\end{definition}

Dua \omterm{def:g10:algebra:equation}{persamaan} disebut \emph{ekuivalen} apabila keduanya mempunyai
penyelesaian yang persis sama. Menambahkan bilangan yang sama pada kedua
ruas, atau mengalikan kedua ruas dengan bilangan \emph{taknol} yang sama,
menghasilkan \omterm{def:g10:algebra:equation}{persamaan} yang ekuivalen.

\begin{example}[Persamaan derajat satu]\label{ex:g10:algebra:firstdegree}
Selesaikan $5x - 7 = 2x + 8$, selangkah demi selangkah:
\begin{align*}
5x - 7 &= 2x + 8 \\
5x - 2x &= 8 + 7 && \text{(tambahkan $7 - 2x$ pada kedua ruas)} \\
3x &= 15 \\
x &= 5 && \text{(bagi kedua ruas dengan $3$).}
\end{align*}
Penyelesaian tunggalnya adalah $5$. Periksa: $5 \times 5 - 7 = 18$ dan
$2 \times 5 + 8 = 18$.
\end{example}

\begin{theorem}[Aturan hasil kali nol]\label{thm:g10:algebra:zeroproduct}
Hasil kali \omterm{def:g10:numbers:sets}{bilangan real} sama dengan nol jika dan hanya jika paling sedikit
satu faktornya nol:
\[
A \times B = 0
\quad\Longleftrightarrow\quad
A = 0 \ \text{ atau } \ B = 0 .
\]
\end{theorem}

\begin{proof}
Jika $A = 0$ atau $B = 0$, jelas hasil kalinya $0$. Sebaliknya, andaikan
$AB = 0$ dan $A \neq 0$. Membagi kedua ruas dengan $A$ memberi $B = 0$.
\end{proof}

\begin{method}[Menyelesaikan dengan memfaktorkan]\label{met:g10:algebra:zeroproduct}
Untuk menyelesaikan \omterm{def:g10:algebra:equation}{persamaan} yang ruas kanannya $0$ dan ruas kirinya dapat
difaktorkan:
\begin{enumerate}
  \item pindahkan semuanya ke kiri sehingga \omterm{def:g10:algebra:equation}{persamaannya} berbunyi $E = 0$;
  \item faktorkan $E$ menjadi hasil kali faktor-faktor derajat satu;
  \item terapkan aturan hasil kali nol lalu selesaikan setiap \omterm{def:g10:algebra:equation}{persamaan} kecil;
  \item kumpulkan semua penyelesaian yang diperoleh.
\end{enumerate}
\end{method}

\begin{example}\label{ex:g10:algebra:zeroproduct}
Selesaikan $(x - 2)(3x + 6) = 0$: entah $x - 2 = 0$, yang memberi $x = 2$,
atau $3x + 6 = 0$, yang memberi $x = -2$. Penyelesaian: $-2$ dan $2$.

Selesaikan $x^2 = 9x$. \emph{Jangan} membagi dengan $x$ (nilainya bisa nol!);
pindahkan semuanya ke satu ruas lalu faktorkan:
\[
x^2 - 9x = 0
\quad\Longleftrightarrow\quad
x(x - 9) = 0
\quad\Longleftrightarrow\quad
x = 0 \ \text{ atau } \ x = 9 .
\]
\end{example}

\begin{example}[Persamaan berbentuk pecahan]\label{ex:g10:algebra:quotient}
Selesaikan $\dfrac{x - 1}{x + 2} = 0$. Sebuah pecahan bernilai nol tepat
ketika pembilangnya nol \emph{dan} penyebutnya tidak. Di sini $x - 1 = 0$
memberi $x = 1$, dan $1 + 2 = 3 \neq 0$: satu-satunya penyelesaian adalah $1$.
Nilai $x = -2$ \emph{terlarang} (pembagian oleh nol) dan harus selalu
dikecualikan lebih dahulu.
\end{example}

\section{Pertidaksamaan dan tabel tanda}

\begin{proposition}[Aturan untuk pertidaksamaan]\label{prop:g10:algebra:ineqrules}
Misalkan $a \leq b$.
\begin{enumerate}
  \item Untuk sebarang $c$: $a + c \leq b + c$ (penjumlahan mempertahankan
        urutan).
  \item Jika $c > 0$: $ac \leq bc$ (mengalikan dengan bilangan positif
        mempertahankan urutan).
  \item Jika $c < 0$: $ac \geq bc$ (mengalikan dengan bilangan negatif
        \emph{membalik} urutan).
\end{enumerate}
\end{proposition}

\begin{proof}
1. $(b + c) - (a + c) = b - a \geq 0$.
2. $bc - ac = (b-a)c$ adalah hasil kali dua bilangan taknegatif, sehingga
$bc \geq ac$.
3. Kini $(b-a)c \leq 0$ sebagai hasil kali bilangan taknegatif dan bilangan
negatif, sehingga $bc \leq ac$.
\end{proof}

\begin{example}\label{ex:g10:algebra:firstineq}
Selesaikan $-2x + 3 > 9$:
\begin{align*}
-2x + 3 &> 9 \\
-2x &> 6 && \text{(kurangi 3)} \\
x &< -3 && \text{(bagi dengan $-2 < 0$: balik arah pertidaksamaan).}
\end{align*}
Himpunan penyelesaiannya adalah $\intoo{-\infty}{-3}$.
\end{example}

\begin{definition}[Tabel tanda]\label{def:g10:algebra:signtable}
Sebuah \emph{tabel tanda}\index{tabel tanda} mencatat, satu baris untuk setiap
faktor, tanda ($+$ atau $-$) setiap faktor sebuah bentuk aljabar ketika $x$
menyusuri $\R$, dengan $0$ pada setiap nilai yang membuat faktor itu lenyap;
baris terakhir memberi tanda hasil kalinya, memakai kaidah tanda kolom demi kolom.
\end{definition}

\begin{example}\label{ex:g10:algebra:signtable}
Tanda dari $P(x) = (x - 2)(3 - x)$. Faktor $x - 2$ lenyap di $2$ dan positif
sesudahnya; faktor $3 - x$ lenyap di $3$ dan positif sebelumnya:
\begin{center}
\renewcommand{\arraystretch}{1.3}
\begin{tabular}{c|ccccc}
$x$ & $-\infty$ \quad & $2$ & & $3$ & \quad $+\infty$ \\
\hline
$x - 2$ & $-$ & $0$ & $+$ & & $+$ \\
$3 - x$ & $+$ & & $+$ & $0$ & $-$ \\
\hline
$P(x)$ & $-$ & $0$ & $+$ & $0$ & $-$ \\
\end{tabular}
\end{center}
Jadi $P(x) > 0$ tepat pada $\intoo{2}{3}$, dan $P(x) \leq 0$ pada
$\intoc{-\infty}{2} \cup \intco{3}{+\infty}$.
\end{example}

\begin{omfigure}
\begin{tikzpicture}
\begin{axis}[omaxis,
  width=8.2cm, height=5.2cm,
  xmin=-0.6, xmax=5.4, ymin=-2.6, ymax=1.4,
  xtick={2,3}, ytick=\empty]
  \addplot[omDef, thick, domain=0.35:4.65] {(x-2)*(3-x)};
  \fill (axis cs:2,0) circle (1.5pt);
  \fill (axis cs:3,0) circle (1.5pt);
  \node[omDef, font=\small] at (axis cs:2.5,0.65) {$P > 0$};
  \node[omThm, font=\small] at (axis cs:0.9,-1.3) {$P < 0$};
  \node[omThm, font=\small] at (axis cs:4.1,-1.3) {$P < 0$};
\end{axis}
\end{tikzpicture}

{\small Grafik $P(x) = (x-2)(3-x)$ membenarkan \omterm{def:g10:algebra:signtable}{tabel tandanya}: positif di
antara kedua akarnya, $2$ dan $3$, serta negatif di luarnya.}
\end{omfigure}

\begin{method}[Menyelesaikan pertidaksamaan dengan tabel tanda]\label{met:g10:algebra:signtable}
Untuk menyelesaikan pertidaksamaan seperti $E(x) \geq 0$ atau $E(x) < 0$:
\begin{enumerate}
  \item pindahkan semuanya ke ruas kiri, sehingga ruas kanannya
        $0$ (jangan pernah mengalikan kedua ruas dengan bentuk aljabar yang
        tandanya belum diketahui);
  \item faktorkan ruas kiri, atau samakan penyebutnya bila ruas itu berupa
        pecahan;
  \item susun \omterm{def:g10:algebra:signtable}{tabel tanda} dengan satu baris untuk setiap faktor (untuk
        pecahan, tandai nilai terlarang dengan garis rangkap pada baris terakhir);
  \item bacalah \omterm{def:g10:numbers:interval}{interval} tempat tanda yang diminta berlaku, sambil memeriksa
        apakah setiap batasnya termasuk.
\end{enumerate}
\end{method}

\begin{example}\label{ex:g10:algebra:quotientineq}
Selesaikan $\dfrac{x + 1}{x - 3} \leq 0$. Pembilangnya lenyap di $-1$,
penyebutnya di $3$ (nilai terlarang). Pecahan itu bernilai negatif ketika
keduanya bertanda berlawanan, yaitu untuk $x$ antara $-1$ dan $3$. Batas
tempat \emph{pembilang} lenyap disertakan, sedangkan nilai terlarang
dikecualikan: himpunan penyelesaiannya adalah $\intco{-1}{3}$.
\end{example}

\section{Latihan}

\begin{exercise}[$\star$]\label{exo:g10:algebra:1}
Jabarkan lalu sederhanakan:
\[
3(2x - 5) - 2(x - 7), \qquad
(x + 4)(2x - 1), \qquad
(3x - 2)^2, \qquad
(5 - x)(5 + x).
\]
\end{exercise}

\begin{exercise}[$\star$]\label{exo:g10:algebra:2}
Faktorkan:
\[
x^2 - 49, \qquad
x^2 + 10x + 25, \qquad
7x^2 - 14x, \qquad
16x^2 - 8x + 1 .
\]
\end{exercise}

\begin{exercise}[$\star$]\label{exo:g10:algebra:3}
Selesaikan:
\[
4x + 9 = x - 3, \qquad
\frac{2x - 1}{3} = 5, \qquad
2(x - 3) = 2x + 1 .
\]
(Salah satunya tidak mempunyai penyelesaian --- jelaskan mengapa.)
\end{exercise}

\begin{exercise}[$\star$]\label{exo:g10:algebra:4}
Selesaikan dengan aturan hasil kali nol:
\[
(x + 5)(2x - 8) = 0, \qquad
x^2 - 16 = 0, \qquad
x^2 + 6x + 9 = 0 .
\]
\end{exercise}

\begin{exercise}[$\star$]\label{exo:g10:algebra:5}
Selesaikan pertidaksamaan berikut, lalu tulis himpunan penyelesaiannya sebagai \omterm{def:g10:numbers:interval}{interval}:
\[
3x - 5 \leq 7, \qquad
-4x + 1 > 9, \qquad
2(x - 1) \geq 5x + 4 .
\]
\end{exercise}

\begin{exercise}[$\star\star$]\label{exo:g10:algebra:6}
Faktorkan $E = (x - 1)(x + 3) - (x - 1)(2x - 5)$, lalu selesaikan $E = 0$.
\end{exercise}

\begin{exercise}[$\star\star$]\label{exo:g10:algebra:7}
Selesaikan $x^2 = 5x$, lalu $(2x + 1)^2 = (x - 4)^2$. (Untuk yang kedua,
pindahkan semuanya ke kiri lalu faktorkan selisih dua kuadratnya.)
\end{exercise}

\begin{exercise}[$\star\star$]\label{exo:g10:algebra:8}
Dengan \omterm{def:g10:algebra:signtable}{tabel tanda}, selesaikan
\[
(x + 2)(4 - x) > 0,
\qquad
\frac{2x - 6}{x + 1} \geq 0 .
\]
\end{exercise}

\begin{exercise}[$\star\star$]\label{exo:g10:algebra:9}
Sebuah layanan siaran daring berbiaya $9$ per bulan, ditambah biaya
pendaftaran sekali bayar sebesar $15$; pesaingnya berbiaya $12$ per bulan
tanpa biaya pendaftaran. Setelah berapa bulan layanan pertama menjadi
pilihan yang lebih murah secara keseluruhan? Susun lalu selesaikan
sebuah pertidaksamaan.
\end{exercise}

\begin{exercise}[$\star\star$]\label{exo:g10:algebra:10}
Selesaikan $\dfrac{x+3}{x-1} = 2$. (Kalikan kedua ruas dengan $x - 1$ setelah
mengecualikan nilai terlarangnya, atau pindahkan semuanya ke kiri dengan
penyebut yang sama.)
\end{exercise}

\begin{exercise}[$\star\star\star$]\label{exo:g10:algebra:11}
Misalkan $n$ \omterm{def:g10:numbers:sets}{bilangan bulat}. Tunjukkan bahwa $(n + 1)^2 - n^2$ selalu ganjil,
dan bahwa selisih kuadrat dua bilangan ganjil berurutan selalu merupakan
kelipatan $8$.
\end{exercise}

\section{Soal: Dari Babilonia ke rasio emas}

\begin{problem}[{Soal akhir pekan --- soal jumlah-dan-hasil-kali,
melengkapkan kuadrat sempurna sebagaimana digambar al-Khwarizmi, dan
persamaan $x^2 = x + 1$}]
\label{pb:g10:algebra:1}
Empat ribu tahun silam, juru tulis Babilonia menyelesaikan soal ini di atas
lempung: \emph{carilah dua bilangan bila diketahui jumlah dan hasil kalinya}.
Kiat mereka --- mencari kedua bilangan secara simetris di sekitar
setengah jumlahnya --- sampai hari ini masih memecahkan setiap \omterm{def:g10:algebra:equation}{persamaan}
kuadrat, dengan nama modernnya, \emph{melengkapkan kuadrat sempurna}. Soal
ini mempelajari kiat itu, menyaksikan al-Khwarizmi benar-benar melengkapkan
sebuah persegi, lalu menerapkan metodenya pada perbandingan yang paling
termasyhur di antara semuanya.

\medskip
\noindent\textbf{Bagian I --- Soal Babilonia.}
\begin{enumerate}
  \item Pemanasan (\cref{thm:g10:algebra:identities}): faktorkan
        $x^2 - 10x + 25$, lalu $9x^2 - 49$, lalu --- dengan memadukan
        sebuah identitas dan selisih dua kuadrat ---
        $x^2 + 6x + 9 - 4$.
  \item Selesaikan $(x + 1)(x + 5) = 0$ dan
        $(3x - 7)(3x + 7) = 0$
        (\cref{thm:g10:algebra:zeroproduct}).
  \item Soal sang juru tulis: dua bilangan berjumlah $10$ dan
        berhasil kali $21$. Mengikuti Babilonia, tulislah keduanya sebagai
        $5 - t$ dan $5 + t$ (simetris di sekitar setengah jumlahnya) lalu
        carilah $t$, kemudian kedua bilangan itu.
  \item Dengan metode yang sama: jumlah $14$, hasil kali $33$.
  \item Resep umumnya: untuk jumlah $s$ dan hasil kali $p$, tunjukkan
        bahwa kedua bilangan itu adalah
        \[
        \frac s2 - t
        \quad\text{dan}\quad
        \frac s2 + t,
        \qquad\text{dengan}\quad
        t^2 = \left(\frac s2\right)^2 - p ,
        \]
        lalu nyatakan syarat pada $s$ dan $p$ agar soal itu
        mempunyai penyelesaian.
\end{enumerate}

\medskip
\noindent\textbf{Bagian II --- Melengkapkan kuadrat sempurna.}
\begin{enumerate}[resume]
  \item Jabarkan $(x - a)(x - b)$ lalu simpulkan: penyelesaian
        $x^2 - sx + p = 0$ tepat adalah pasangan bilangan yang
        berjumlah $s$ dan berhasil kali $p$. (Para juru tulis
        menyelesaikan \omterm{def:g10:algebra:equation}{persamaan} kuadrat tanpa menyadarinya.)
  \item Periksa bahwa $x^2 - 10x + 21 = (x - 5)^2 - 4$, lalu selesaikan
        $x^2 - 10x + 21 = 0$ dengan \omterm{def:g10:algebra:expand}{memfaktorkan} ruas kanannya sebagai
        selisih dua kuadrat. Bandingkan dengan pertanyaan 3.
  \item Selesaikan $x^2 + 6x - 7 = 0$ dengan melengkapkan kuadrat sempurna.
  \item Lengkapkan kuadrat sempurna pada $x^2 + 4x + 5$, lalu jelaskan
        mengapa \omterm{def:g10:algebra:equation}{persamaan} $x^2 + 4x + 5 = 0$ \emph{tidak} mempunyai
        penyelesaian real. Nyatakan kriteria umumnya: setelah ditulis
        sebagai $(x + h)^2 + k$, kapan sebuah bentuk kuadrat lenyap?
  \item Namanya harfiah. Di Bagdad pada abad kesembilan,
        al-Khwarizmi menggambar $x^2 + 6x = 7$ sebagai persegi bersisi $x$
        dengan dua persegi panjang $3 \times x$ yang direkatkan pada dua
        sisinya, lalu \emph{melengkapkan persegi itu} dengan pojok
        $3 \times 3$ yang hilang. Gambarlah bangunnya dan jelaskan bagaimana
        gambar itu mengubah \omterm{def:g10:algebra:equation}{persamaannya} menjadi $(x + 3)^2 = 16$ --- lalu
        selesaikan.
\end{enumerate}

\medskip
\noindent\textbf{Bagian III --- Tanda, tabel dan dua janji
lama.}
\begin{enumerate}[resume]
  \item Susun \omterm{def:g10:algebra:signtable}{tabel tanda} dari $(x - 3)(x + 2)$
        (\cref{def:g10:algebra:signtable}) lalu selesaikan
        $(x - 3)(x + 2) < 0$.
  \item Selesaikan $x^2 \leq 10x - 21$ (pindahkan semuanya ke satu ruas
        lalu pakai kembali pertanyaan 7).
  \item Selesaikan $\dfrac{x - 1}{x + 2} \geq 0$ dengan \omterm{def:g10:algebra:signtable}{tabel tanda},
        sambil memperhatikan nilai terlarangnya.
  \item Sebuah janji lama ditepati: bagi petani pada soal pagar Ratu Dido
        di jilid sebelumnya, dengan pagar sepanjang $20$ m, luas
        kandang berukuran $x \times (10 - x)$ memenuhi
        \[
        A(x) = x(10 - x) = 25 - (x - 5)^2 .
        \]
        Periksa identitas itu, lalu bacalah darinya sebuah \emph{bukti}
        bagi apa yang di jilid sebelumnya hanya dapat diduga: luas
        maksimumnya, dan satu-satunya $x$ yang mencapainya.
  \item Janji lain: buktikan bahwa $x + \dfrac1x \geq 2$ untuk setiap
        $x > 0$, dengan kesamaan hanya di $x = 1$. (Kalikan dengan
        $x$, kumpulkan sukunya, kenali sebuah kuadrat --- ketaksamaan
        rataan aritmetika--geometri pada soal akhir pekan tentang relasi
        Euclid di jilid sebelumnya dalam wujud baru.)
\end{enumerate}

\medskip
\noindent\textbf{Bagian IV --- Potongan emas.}
Potonglah sebuah ruas garis sepanjang $1$ menjadi bagian besar $x$ dan
bagian kecil $1 - x$ sedemikian sehingga
\emph{seluruhnya berbanding bagian besar sama dengan bagian besar berbanding bagian kecil}:
$\frac{1}{x} = \frac{x}{1 - x}$.
\begin{enumerate}[resume]
  \item Tunjukkan bahwa syarat itu berbunyi $x^2 + x - 1 = 0$.
  \item Selesaikan dengan melengkapkan kuadrat sempurna, lalu simpan akar
        yang berupa panjang: $x = \frac{\sqrt5 - 1}{2} \approx 0.618$.
  \item \emph{Rasio emas} adalah
        $\varphi = \frac1x = \frac{2}{\sqrt5 - 1}$. Rasionalkan
        penyebutnya (kalikan pembilang dan penyebut dengan
        $\sqrt5 + 1$, sebuah identitas sedang bekerja) lalu tunjukkan
        $\varphi = \frac{\sqrt5 + 1}{2} \approx 1.618$.
  \item Periksa \emph{secara eksak} kedua identitas ajaib
        $\varphi$:
        \[
        \varphi^2 = \varphi + 1,
        \qquad
        \frac{1}{\varphi} = \varphi - 1 .
        \]
  \item Penutup: jelaskan mengapa $\varphi$ bersifat \omterm{ex:g10:numbers:classify}{irasional}
        (soal akhir pekan tentang sifat \omterm{ex:g10:numbers:classify}{irasional} di jilid sebelumnya
        memberi fakta kuncinya tentang $\sqrt5$), lalu hitunglah rasio
        $\frac53$, $\frac85$, $\frac{13}{8}$, $\frac{21}{13}$ dari
        bilangan Hemachandra--Fibonacci yang berurutan (soal akhir pekan
        tentang pencacahan irama di jilid sebelumnya) sampai tiga angka
        desimal. Apa yang tampaknya mereka kejar? (Bukti bahwa mereka
        benar-benar menangkapnya menjadi bagian bab tentang barisan.)

\end{enumerate}
\end{problem}
